\documentclass[a4paper]{article}
\usepackage[margin=2.5cm]{geometry}
\usepackage{graphicx} % Required for inserting images
\usepackage[utf8]{inputenc}
\usepackage{lmodern}
\usepackage[T1]{fontenc} % portable Lithuanian font encoding
\usepackage[provide=*]{babel}
\usepackage{array}
\usepackage{tabularx}
\usepackage{booktabs}
\usepackage{longtable}
\usepackage{makecell}
\usepackage{float}
\usepackage{pgfgantt} % Ganto diagrama (tvarkaraštis)
\usepackage{pdflscape} % gulsčias puslapis Ganto diagramai
\usepackage{caption}
% babel lt: „1 lentelė“, „1 pav.“; lentelės pavadinimas virš lentelės, paveikslo – po paveikslu
\captionsetup[table]{labelsep=period}
\captionsetup[figure]{labelsep=space}

\babelprovide[main, import]{lithuanian}




\begin{document}

\begin{titlepage}
    \centering
    
    {\Large \textbf{Vilniaus Universitetas}}\\[0.2cm]
    {\small Matematikos ir informatikos fakultetas}\\[0.1cm]
    {\small Informatikos katedra}\\[4cm]

    {\huge \textbf{Reikalavimų specifikacija}}\\[4cm]

    \begin{raggedleft}
    \begin{minipage}{0.55\textwidth}
        \textbf{Atliko:}\\
        Tomas Baublys, Tadas Plančiūnas, Goda Gutparakytė, Marijus Plančiūnas, Ugnė Skučaitė\\[0.8cm]
    \end{minipage}
    \end{raggedleft}

    \vfill

    Vilnius \\ \the\year

\end{titlepage}

\pagebreak

\tableofcontents
\pagebreak

% Funkciniai reikalavimai - Ką sistema turi atlikti. Kiekvienas reikalavimas turi unikalų ID, aiškų aprašymą ir patikrinamą rezultatą. GODA
% Nefunkciniai reikalavimai - Sistemos kokybės ir veikimo ribojimai: našumas, saugumas, patikimumas, naudojamumas, suderinamumas. Formuluojame išmatuojamai. UGNĖ
% Atsekamumas su poreikiais - Lentelė, susiejanti poreikio ID su funkcinių ir nefunkcinių reikalavimų ID. Papildomai galima nurodyti susijusį panaudojimo atvejį. GODA

\section{Funkciniai reikalavimai}

Šiame skyriuje pateikiami visi sistemai keliami funkciniai reikalavimai, apibrėžiantys, ką sistema turi atlikti. Kiekvienas reikalavimas turi unikalų identifikatorių (\textbf{FR-XX}), aiškų aprašymą, patikrinamą rezultatą (testavimo kriterijų) ir yra tiesiogiai susietas su užsakovo poreikiais bei sistemos dalimis. Smulkesnis sistemos funkcijų išskaidymas pagal kodus (P0.1--P7.3) pateikiamas projekto plano A priedo 18 lentelėje.

Panaudojimo atvejų (PA-01--PA-18) identifikatoriai paimti iš Analizės modelio skyriaus Panaudojimo atvejų registro. Prie kiekvieno reikalavimo nurodyti tie panaudojimo atvejai, kurių šaltinio funkcijos (P0--P7 kodai) sutampa su reikalavimą atitinkančiomis funkcijomis. Vienas panaudojimo atvejis gali apimti kelias susijusias funkcijas, todėl tas pats PA gali būti nurodytas prie kelių reikalavimų. Bendras poreikių, reikalavimų ir panaudojimo atvejų ryšys pateikiamas skyriuje „Atsekamumas su poreikiais“.

\subsection{Sistemos dalių struktūra}
Lentelėje pateikiamas pagrindinių sistemos dalių (kodų) pasiskirstymas, pagal kurį grupuojami reikalavimai.

\begin{table}[h!]
\centering
\begin{tabular}{|l|l|}
\hline
\textbf{Kodas} & \textbf{Sistemos dalis} \\ \hline
P0 & Paskyros, administravimas ir sporto erdvės \\ \hline
P1 & Užsiėmimų programų valdymas \\ \hline
P2 & Užsiėmimų valdymas \\ \hline
P3 & Dalyvių valdymas \\ \hline
P4 & Asmeninių programų valdymas \\ \hline
P5 & Vaizdinio turinio pateikimas \\ \hline
P6 & Garso įranga \\ \hline
P7 & Duomenų ir istorijos valdymas \\ \hline
\end{tabular}
\caption{Pagrindinės sistemos dalys (P0--P7)}
\label{tab:sistemos_dalys}
\end{table}

\subsection{P0. Paskyros, administravimas ir sporto erdvės}

\vspace{0.3cm}
\noindent\textbf{FR-01: Dalyvio registracija sistemoje}
\begin{itemize}
    \item \textbf{Užsakovo poreikis ir funkcijos:} F1 (atitinka funkcijos kodą P0.4 iš projekto plano 18 lentelės), panaudojimo atvejis PA-01.
    \item \textbf{Aprašymas:} Sistema turi leisti naujiems dalyviams užsiregistruoti sistemoje per žiniatinklio sąsają, suvedant privalomus asmens duomenis.
    \item \textbf{Patikrinamas rezultatas:} Sukuriama nauja paskyra su dalyvio vaidmeniu, o sėkmingai užsiregistravęs naudotojas gali prisijungti.
\end{itemize}

\vspace{0.3cm}
\noindent\textbf{FR-02: Naudotojo autorizacija (prisijungimas)}
\begin{itemize}
    \item \textbf{Užsakovo poreikis ir funkcijos:} F1 (atitinka funkcijos kodą P0.5 iš projekto plano 18 lentelės), panaudojimo atvejis PA-01.
    \item \textbf{Aprašymas:} Sistema turi leisti registruotiems naudotojams prisijungti prie savo paskyros suvedant el. pašto adresą ir slaptažodį.
    \item \textbf{Patikrinamas rezultatas:} Įvedus teisingus duomenis, naudotojas nukreipiamas į savo vaidmens pradinį puslapį; įvedus neteisingus – parodomas pranešimas.
\end{itemize}

\vspace{0.3cm}
\noindent\textbf{FR-03: Užmiršto slaptažodžio atkūrimas}
\begin{itemize}
    \item \textbf{Užsakovo poreikis ir funkcijos:} F1 (atitinka funkcijos kodą P0.6 iš projekto plano 18 lentelės), panaudojimo atvejis PA-01.
    \item \textbf{Aprašymas:} Sistema turi suteikti galimybę naudotojui inicijuoti slaptažodžio atkūrimą el. paštu išsiunčiant vienkartinę saugią nuorodą.
    \item \textbf{Patikrinamas rezultatas:} Naudotojas gauna el. laišką, paspaudžia nuorodą ir sėkmingai atnaujina savo slaptažodį.
\end{itemize}

\vspace{0.3cm}
\noindent\textbf{FR-04: Naudotojo asmeninio profilio peržiūra ir redagavimas}
\begin{itemize}
    \item \textbf{Užsakovo poreikis ir funkcijos:} F1, F2 (atitinka funkcijos kodus P0.1, P0.7, P0.8 iš projekto plano 18 lentelės), panaudojimo atvejis PA-01.
    \item \textbf{Aprašymas:} Prisijungęs naudotojas turi turėti galimybę peržiūrėti bei atnaujinti savo asmeninio profilio duomenis (kontaktus, slaptažodį).
    \item \textbf{Patikrinamas rezultatas:} Atlikti pakeitimai išsaugomi sistemoje ir atsispindi naudotojo profilio lange.
\end{itemize}

\vspace{0.3cm}
\noindent\textbf{FR-05: Vaidmenimis pagrįsta prieigos kontrolė (RBAC)}
\begin{itemize}
    \item \textbf{Užsakovo poreikis ir funkcijos:} F2 (atitinka funkcijos kodą P0.11 iš projekto plano 18 lentelės), panaudojimo atvejai PA-01, PA-02.
    \item \textbf{Aprašymas:} Sistema turi užtikrinti trijų vaidmenų (Dalyvis, Treneris, Administratorius) teisių suteikimą bei griežtai riboti prieigą prie neleistinų funkcijų.
    \item \textbf{Patikrinamas rezultatas:} Bandant pasiekti kito vaidmens funkciją, sistema grąžina prieigos klaidos pranešimą (pvz., Dalyvis nemato kitų dalyvių asmens duomenų).
\end{itemize}

\vspace{0.3cm}
\noindent\textbf{FR-06: Naudotojų paskyrų sukūrimas ir redagavimas administratoriaus sąsajoje}
\begin{itemize}
    \item \textbf{Užsakovo poreikis ir funkcijos:} F2, F14 (atitinka funkcijos kodus P0.9, P0.10 iš projekto plano 18 lentelės), panaudojimo atvejis PA-02.
    \item \textbf{Aprašymas:} Administratorius turi turėti galimybę tiesiogiai sukurti naują naudotoją bei redaguoti esamų naudotojų paskyros duomenis ir vaidmenis.
    \item \textbf{Patikrinamas rezultatas:} Administratoriaus sukurta paskyra gali sėkmingai prisijungti, o pakeisti duomenys išsaugomi duomenų bazėje.
\end{itemize}

\vspace{0.3cm}
\noindent\textbf{FR-07: Paskyros deaktyvavimas}
\begin{itemize}
    \item \textbf{Užsakovo poreikis ir funkcijos:} F2, F14 (atitinka funkcijos kodą P0.12 iš projekto plano 18 lentelės), panaudojimo atvejis PA-02.
    \item \textbf{Aprašymas:} Administratorius turi turėti galimybę deaktyvuoti naudotojo paskyrą. Deaktyvuojant dalyvį, jo būsimos rezervacijos automatiškai atšaukiamos; trenerį galima deaktyvuoti tik jei jis neturi būsimų užsiėmimų.
    \item \textbf{Patikrinamas rezultatas:} Deaktyvuotas naudotojas nebegali prisijungti, jo būsimos rezervacijos anuliuojamos; bandant deaktyvuoti įsipareigojimų turintį trenerį, sistema atmeta veiksmą.
\end{itemize}

\vspace{0.3cm}
\noindent\textbf{FR-08: Naudotojų sąrašo ir detalių peržiūra}
\begin{itemize}
    \item \textbf{Užsakovo poreikis ir funkcijos:} F14 (atitinka funkcijos kodus P0.13, P0.14 iš projekto plano 18 lentelės), panaudojimo atvejis PA-02.
    \item \textbf{Aprašymas:} Administratorius turi turėti galimybę peržiūrėti visų sistemos naudotojų sąrašą, filtruoti jį pagal vaidmenį bei atidaryti konkretaus naudotojo detalią paskyros kortelę.
    \item \textbf{Patikrinamas rezultatas:} Sąraše teisingai pritaikomas filtras pagal vaidmenį, o atidarius kortelę parodoma pilna informacija apie naudotoją.
\end{itemize}

\vspace{0.3cm}
\noindent\textbf{FR-09: Sporto erdvių registras}
\begin{itemize}
    \item \textbf{Užsakovo poreikis ir funkcijos:} F14, F15 (atitinka funkcijos kodus P0.2, P0.15, P0.16, P0.17, P0.18 iš projekto plano 18 lentelės), panaudojimo atvejis PA-03.
    \item \textbf{Aprašymas:} Administratorius turi turėti galimybę pridėti naujas sporto erdves, redaguoti jų profilį, peržiūrėti sąrašą bei šalinti erdves (tikrinant, ar jose nėra suplanuotų užsiėmimų).
    \item \textbf{Patikrinamas rezultatas:} Erdvės duomenys atnaujinami sąrašuose; bandant pašalinti erdvę su aktyviais užsiėmimais, trynimas užblokuojamas.
\end{itemize}

\vspace{0.3cm}
\noindent\textbf{FR-10: Techninės įrangos valdymo ir nustatymų registras}
\begin{itemize}
    \item \textbf{Užsakovo poreikis ir funkcijos:} F14, F15 (atitinka funkcijos kodus P0.3, P0.19, P0.20, P0.21, P0.22, P0.23 iš projekto plano 18 lentelės), panaudojimo atvejis PA-03.
    \item \textbf{Aprašymas:} Administratorius turi galimybę priskirti techninę įrangą (grotuvus, ekranus, garso siųstuvus) prie konkrečių erdvių, keisti ir peržiūrėti įrenginių vaizdo/garso parametrus bei juos pašalinti.
    \item \textbf{Patikrinamas rezultatas:} Įrangos parametrai išsaugomi duomenų bazėje ir atvaizduojami atitinkamos erdvės įrangos nustatymų lange.
\end{itemize}

\subsection{P1. Užsiėmimų programų valdymas}

\vspace{0.3cm}
\noindent\textbf{FR-11: Mokomosios vaizdo ir garso medžiagos valdymas}
\begin{itemize}
    \item \textbf{Užsakovo poreikis ir funkcijos:} F6, F14 (atitinka funkcijos kodus P1.2, P1.7, P1.8, P1.9, P1.18, P1.19 iš projekto plano 18 lentelės), panaudojimo atvejis PA-04.
    \item \textbf{Aprašymas:} Administratorius turi turėti galimybę įkelti medijos failus (vaizdo/garso įrašus) su metaduomenimis, juos perklausyti/peržiūrėti, redaguoti metaduomenis bei šalinti failus (tikrinant panaudojimą pratimuose).
    \item \textbf{Patikrinamas rezultatas:} Įkeltas failas išsaugomas serveryje; bandant ištrinti pratimui priskirtą mediją, sistema parodo perspėjimą ir trynimą atmeta.
\end{itemize}

\vspace{0.3cm}
\noindent\textbf{FR-12: Pratimų katalogo valdymas ir paieška}
\begin{itemize}
    \item \textbf{Užsakovo poreikis ir funkcijos:} F6, F14 (atitinka funkcijos kodus P1.1, P1.4, P1.5, P1.6, P1.16, P1.17 iš projekto plano 18 lentelės), panaudojimo atvejis PA-04.
    \item \textbf{Aprašymas:} Administratorius turi turėti galimybę kurti, redaguoti ir šalinti pratimus kataloge bei ieškoti pratimų su medijos demonstracija. Šalinant tikrinama, ar pratimas nenaudojamas programose.
    \item \textbf{Patikrinamas rezultatas:} Naujas pratimas rodomas kataloge; paieška grąžina atitinkančius įrašus; naudojamo pratimo šalinimas užblokuojamas.
\end{itemize}

\vspace{0.3cm}
\noindent\textbf{FR-13: Užsiėmimo programos kūrimas ir peržiūra}
\begin{itemize}
    \item \textbf{Užsakovo poreikis ir funkcijos:} F5, F7 (atitinka funkcijos kodus P1.3, P1.10, P1.20, P1.21 iš projekto plano 18 lentelės), panaudojimo atvejis PA-05.
    \item \textbf{Aprašymas:} Treneris turi turėti galimybę sukurti naują užsiėmimo programą bei peržiūrėti programų sąrašą ir išsamią programos struktūrą (etapus, pratimus, mediją).
    \item \textbf{Patikrinamas rezultatas:} Sukurta programa išsaugoma sistemoje ir tampa matoma programų sąraše su visais priskirtais elementais.
\end{itemize}

\vspace{0.3cm}
\noindent\textbf{FR-14: Programos etapų valdymas}
\begin{itemize}
    \item \textbf{Užsakovo poreikis ir funkcijos:} F5, F7 (atitinka funkcijos kodus P1.13, P1.14, P1.15 iš projekto plano 18 lentelės), panaudojimo atvejis PA-05.
    \item \textbf{Aprašymas:} Treneris turi galimybę prie programos pridėti etapus, nustatant trukmę, tekstines instrukcijas, pratimus bei mediją, keisti etapų eiliškumą arba etapus šalinti.
    \item \textbf{Patikrinamas rezultatas:} Pakeitus etapų eiliškumą ar parametrus, programos peržiūros lange iškart rodoma atnaujinta struktūra.
\end{itemize}

\vspace{0.3cm}
\noindent\textbf{FR-15: Programų versijavimas, archyvavimas ir istorija}
\begin{itemize}
    \item \textbf{Užsakovo poreikis ir funkcijos:} F5, F7 (atitinka funkcijos kodus P1.11, P1.12, P1.22 iš projekto plano 18 lentelės), panaudojimo atvejis PA-05.
    \item \textbf{Aprašymas:} Redaguojant jau esamą programą, sistema turi automatiškai sukurti naują versiją, išsaugodama versijų istoriją su autoriumi, bei leisti nebenaudojamas programas suarchyvuoti.
    \item \textbf{Patikrinamas rezultatas:} Išsaugojus pakeitimus, sugeneruojamas naujas versijos numeris, o versijų istorijos lange galima peržiūrėti senesnes versijas.
\end{itemize}

\subsection{P2. Užsiėmimų valdymas}

\vspace{0.3cm}
\noindent\textbf{FR-16: Grupinių ir individualių užsiėmimų planavimas}
\begin{itemize}
    \item \textbf{Užsakovo poreikis ir funkcijos:} F3 (atitinka funkcijos kodus P2.1, P2.2 iš projekto plano 18 lentelės), panaudojimo atvejis PA-06.
    \item \textbf{Aprašymas:} Treneris arba administratorius turi turėti galimybę suplanuoti grupinį arba individualų užsiėmimą tvarkaraštyje, nurodant datą, laiką, erdvę, programą bei vietų skaičių (individualiam – 1 vieta).
    \item \textbf{Patikrinamas rezultatas:} Suplanuotas užsiėmimas atsiranda tvarkaraštyje. Bandant suplanuoti sutampančio laiko užsiėmimą tam pačiam treneriui, sistema veiksmą atmeta.
\end{itemize}

\vspace{0.3cm}
\noindent\textbf{FR-17: Virtualaus trenerio užsiėmimų planavimas}
\begin{itemize}
    \item \textbf{Užsakovo poreikis ir funkcijos:} F17 (atitinka funkcijos kodą P2.3 iš projekto plano 18 lentelės), panaudojimo atvejis PA-07.
    \item \textbf{Aprašymas:} Tik administratorius turi turėti teisę suplanuoti virtualaus trenerio užsiėmimą tvarkaraštyje pasirinktai erdvei ir programai.
    \item \textbf{Patikrinamas rezultatas:} Suplanuotas virtualaus trenerio užsiėmimas rodomas tvarkaraštyje su atitinkama žyma; trenerio paskyra tam neturi teisių.
\end{itemize}

\vspace{0.3cm}
\noindent\textbf{FR-18: Užsiėmimo duomenų keitimas ir informavimas}
\begin{itemize}
    \item \textbf{Užsakovo poreikis ir funkcijos:} F3 (atitinka funkcijos kodus P2.4, P2.13 iš projekto plano 18 lentelės), panaudojimo atvejis PA-08.
    \item \textbf{Aprašymas:} Treneris arba administratorius turi turėti galimybę modifikuoti suplanuoto užsiėmimo parametrus (laiką, erdvę, vietų skaičių), automatiškai išsiunčiant pranešimus užsiregistravusiems dalyviams.
    \item \textbf{Patikrinamas rezultatas:} Pakeitus užsiėmimo laiką, registruoti dalyviai gauna informacinius pranešimus, o tvarkaraštyje atsispindi nauji duomenys.
\end{itemize}

\vspace{0.3cm}
\noindent\textbf{FR-19: Užsiėmimo atšaukimas ir rezervacijų anuliavimas}
\begin{itemize}
    \item \textbf{Užsakovo poreikis ir funkcijos:} F3 (atitinka funkcijos kodą P2.5 iš projekto plano 18 lentelės), panaudojimo atvejis PA-08.
    \item \textbf{Aprašymas:} Treneris arba administratorius turi galimybę atšaukti suplanuotą užsiėmimą; atšaukus automatiškai anuliuojamos visos dalyvių rezervacijos ir jiems išsiunčiami pranešimai.
    \item \textbf{Patikrinamas rezultatas:} Užsiėmimo statusas pasikeičia į „atšauktas“, visos susijusios rezervacijos anuliuojamos, o dalyviai gauna pranešimus.
\end{itemize}

\vspace{0.3cm}
\noindent\textbf{FR-20: Užsiėmimo eigos paleidimas ir užbaigimas}
\begin{itemize}
    \item \textbf{Užsakovo poreikis ir funkcijos:} F8 (atitinka funkcijos kodus P2.6, P2.10 iš projekto plano 18 lentelės), panaudojimo atvejis PA-15.
    \item \textbf{Aprašymas:} Treneris per valdymo pultą turi galimybę paleisti suplanuotą užsiėmimą arba jį rankiniu būdu užbaigti.
    \item \textbf{Patikrinamas rezultatas:} Paspaudus „Pradėti“, vietiniam grotuvui perduodama paleidimo komanda; paspaudus „Baigti“, užsiėmimo statusas pakeičiamas į „įvykęs“.
\end{itemize}

\vspace{0.3cm}
\noindent\textbf{FR-21: Užsiėmimo eigos pristabdymas ir tęsimas}
\begin{itemize}
    \item \textbf{Užsakovo poreikis ir funkcijos:} F8 (atitinka funkcijos kodą P2.7 iš projekto plano 18 lentelės), panaudojimo atvejis PA-15.
    \item \textbf{Aprašymas:} Treneris per valdymo pultą užsiėmimo metu turi galimybę laikinai sustabdyti (pauzė) ir vėl pratęsti programos vykdymą.
    \item \textbf{Patikrinamas rezultatas:} Paspaudus pauzę, vaizdo/garso atkūrimas ir laikmatis sustoja; paspaudus tęsti – atnaujinamas nuo tos pačios vietos.
\end{itemize}

\vspace{0.3cm}
\noindent\textbf{FR-22: Etapų praleidimas ir eiliškumo keitimas realiuoju laiku}
\begin{itemize}
    \item \textbf{Užsakovo poreikis ir funkcijos:} F8 (atitinka funkcijos kodą P2.8 iš projekto plano 18 lentelės), panaudojimo atvejis PA-15.
    \item \textbf{Aprašymas:} Treneris užsiėmimo metu per valdymo pultą turi galimybę praleisti esamą etapą ir iškart pereiti prie kito pasirinkto etapo.
    \item \textbf{Patikrinamas rezultatas:} Grotuve iškart sustabdomas dabartinis etapas ir pradedamas rodyti bei transliuoti pasirinktas naujas etapas.
\end{itemize}

\vspace{0.3cm}
\noindent\textbf{FR-23: Etapo trukmės koregavimas vykdymo metu}
\begin{itemize}
    \item \textbf{Užsakovo poreikis ir funkcijos:} F8 (atitinka funkcijos kodą P2.9 iš projekto plano 18 lentelės), panaudojimo atvejis PA-15.
    \item \textbf{Aprašymas:} Treneris turi turėti galimybę valdymo pulte realiu laiku pratęsti arba sutrumpinti vykstančio etapo laiko skaičiavimą.
    \item \textbf{Patikrinamas rezultatas:} Ekrane ir pultelyje atvaizduojamas atnaujintas likusio etapo laiko skaičiavimas.
\end{itemize}

\vspace{0.3cm}
\noindent\textbf{FR-24: Valdymo pulto būsenos atvaizdavimas}
\begin{itemize}
    \item \textbf{Užsakovo poreikis ir funkcijos:} F8 (atitinka funkcijos kodus P2.11, P2.12 iš projekto plano 18 lentelės), panaudojimo atvejis PA-15.
    \item \textbf{Aprašymas:} Valdymo pulto sąsaja turi rodyti dabartinį ir sekantį etapus, likusį laiką, užsiėmimo būseną bei erdvių užimtumą.
    \item \textbf{Patikrinamas rezultatas:} Trenerio ekrane realiuoju laiku atnaujinami visi vykstančio užsiėmimo parametrai.
\end{itemize}

\vspace{0.3cm}
\noindent\textbf{FR-25: Virtualaus trenerio autonominis vykdymas}
\begin{itemize}
    \item \textbf{Užsakovo poreikis ir funkcijos:} F17 (atitinka funkcijos kodus P2.3, P2.6 iš projekto plano 18 lentelės), panaudojimo atvejis PA-07, PA-15.
    \item \textbf{Aprašymas:} Prisijungusiam treneriui valdymo pulte paleidus virtualaus trenerio užsiėmimą, vietinis grotuvas turi autonomiškai atlikti visus programos etapus be tolimesnio žmogaus įsikišimo.
    \item \textbf{Patikrinamas rezultatas:} Programa nuosekliai pereina per visus etapus, ekrane rodydama vaizdą ir transliuodama garsą iki pat užsiėmimo pabaigos.
\end{itemize}

\subsection{P3. Dalyvių valdymas}

\vspace{0.3cm}
\noindent\textbf{FR-26: Tvarkaraščio ir laisvų vietų peržiūra}
\begin{itemize}
    \item \textbf{Užsakovo poreikis ir funkcijos:} F4 (atitinka funkcijos kodą P3.2 iš projekto plano 18 lentelės), panaudojimo atvejis PA-09.
    \item \textbf{Aprašymas:} Dalyvis turi galimybę peržiūrėti užsiėmimų tvarkaraštį su dinaminiu laisvų vietų skaičiumi.
    \item \textbf{Patikrinamas rezultatas:} Tvarkaraštyje teisingai rodomas likusių laisvų vietų skaičius kiekviename užsiėmime.
\end{itemize}

\vspace{0.3cm}
\noindent\textbf{FR-27: Vietos rezervavimas užsiėmime}
\begin{itemize}
    \item \textbf{Užsakovo poreikis ir funkcijos:} F4 (atitinka funkcijos kodą P3.3 iš projekto plano 18 lentelės), panaudojimo atvejis PA-09.
    \item \textbf{Aprašymas:} Dalyvis turi turėti galimybę rezervuoti vietą pasirinktame grupiniame užsiėmime arba rezervuoti individualios pamokos laiką.
    \item \textbf{Patikrinamas rezultatas:} Rezervavus vietą, laisvų vietų skaičius sumažėja 1, o įrašas atsiranda dalyvio rezervacijose. Rezervuoti neleidžiama, jei vietų nėra arba paskyra deaktyvuota.
\end{itemize}

\vspace{0.3cm}
\noindent\textbf{FR-28: Rezervacijos atšaukimas dalyvio iniciatyva}
\begin{itemize}
    \item \textbf{Užsakovo poreikis ir funkcijos:} F4 (atitinka funkcijos kodą P3.4 iš projekto plano 18 lentelės), panaudojimo atvejis PA-10.
    \item \textbf{Aprašymas:} Dalyvis turi turėti galimybę savarankiškai atšaukti savo aktyvią būsimo užsiėmimo rezervaciją iki užsiėmimo pradžios.
    \item \textbf{Patikrinamas rezultatas:} Atšaukus rezervaciją, jos statusas pakeičiamas į „atšaukta“, o laisvų vietų skaičius užsiėmime padidėja vienetu.
\end{itemize}

\vspace{0.3cm}
\noindent\textbf{FR-29: Asmeninės rezervacijų istorijos peržiūra}
\begin{itemize}
    \item \textbf{Užsakovo poreikis ir funkcijos:} F4 (atitinka funkcijos kodą P3.8 iš projekto plano 18 lentelės), panaudojimo atvejis PA-09.
    \item \textbf{Aprašymas:} Dalyvis turi galimybę peržiūrėti savo esamų (aktyvių) ir praėjusių rezervacijų sąrašą.
    \item \textbf{Patikrinamas rezultatas:} Sąraše aiškiai atskiriamos būsimos ir praėjusios rezervacijos su atitinkamais statusais.
\end{itemize}

\vspace{0.3cm}
\noindent\textbf{FR-30: Dalyvio priskyrimas prie užsiėmimo rankiniu būdu}
\begin{itemize}
    \item \textbf{Užsakovo poreikis ir funkcijos:} F12 (atitinka funkcijos kodą P3.5 iš projekto plano 18 lentelės), panaudojimo atvejis PA-11.
    \item \textbf{Aprašymas:} Treneris (savo vedamiems ir virtualaus trenerio užsiėmimams) bei administratorius (visiems užsiėmimams) turi turėti galimybę rankiniu būdu priskirti dalyvį prie užsiėmimo. Individualiam užsiėmimui leidžiama priskirti tik 1 dalyvį.
    \item \textbf{Patikrinamas rezultatas:} Priskirtas dalyvis atsiranda užsiėmimo sąraše; bandant priskirti antrą dalyvį individualiam užsiėmimui, sistema veiksmą atmeta.
\end{itemize}

\vspace{0.3cm}
\noindent\textbf{FR-31: Dalyvio pašalinimas iš užsiėmimo sąrašo}
\begin{itemize}
    \item \textbf{Užsakovo poreikis ir funkcijos:} F12 (atitinka funkcijos kodą P3.6 iš projekto plano 18 lentelės), panaudojimo atvejis PA-11.
    \item \textbf{Aprašymas:} Treneris arba administratorius turi turėti galimybę pašalinti dalyvį iš užsiėmimo sąrašo.
    \item \textbf{Patikrinamas rezultatas:} Pašalinto dalyvio rezervacijos būsena pakeičiama į „pašalinta“, o laisvų vietų skaičius atitinkamai padidėja.
\end{itemize}

\vspace{0.3cm}
\noindent\textbf{FR-32: Užsiėmimo dalyvių sąrašo ir kortelių peržiūra}
\begin{itemize}
    \item \textbf{Užsakovo poreikis ir funkcijos:} F12 (atitinka funkcijos kodus P3.9, P3.12, P3.13 iš projekto plano 18 lentelės), panaudojimo atvejis PA-11.
    \item \textbf{Aprašymas:} Treneris turi turėti galimybę peržiūrėti savo užsiėmimų dalyvių sąrašą, ieškoti dalyvio pagal vardą/pavardę bei atidaryti jo informacinę kortelę.
    \item \textbf{Patikrinamas rezultatas:} Treneris mato tik savo vedamų bei virtualaus trenerio užsiėmimų dalyvius ir gali peržiūrėti jų duomenis.
\end{itemize}

\vspace{0.3cm}
\noindent\textbf{FR-33: Faktinio lankomumo žymėjimas ir taisymas}
\begin{itemize}
    \item \textbf{Užsakovo poreikis ir funkcijos:} F13 (atitinka funkcijos kodą P3.7 iš projekto plano 18 lentelės), panaudojimo atvejis PA-12.
    \item \textbf{Aprašymas:} Treneris užsiėmimo metu arba po jo turi galimybę pažymėti arba pakoreguoti kiekvieno užsiregistravusio dalyvio faktinį lankomumą (dalyvavo / nedalyvavo).
    \item \textbf{Patikrinamas rezultatas:} Lankomumo žymos išsaugomos sistemoje ir atsispindi lankomumo ataskaitose.
\end{itemize}

\vspace{0.3cm}
\noindent\textbf{FR-34: Lankomumo istorijos peržiūra}
\begin{itemize}
    \item \textbf{Užsakovo poreikis ir funkcijos:} F13 (atitinka funkcijos kodą P3.10 iš projekto plano 18 lentelės), panaudojimo atvejis PA-12.
    \item \textbf{Aprašymas:} Dalyvis (savo) ir treneris (savo klientų) turi turėti galimybę peržiūrėti faktinio lankomumo istoriją.
    \item \textbf{Patikrinamas rezultatas:} Sąsajoje pateikiami tikslūs praėjusių užsiėmimų lankomumo įrašai.
\end{itemize}

\vspace{0.3cm}
\noindent\textbf{FR-35: Automatiniai priminimai apie užsiėmimą}
\begin{itemize}
    \item \textbf{Užsakovo poreikis ir funkcijos:} F4 (atitinka funkcijos kodą P3.11 iš projekto plano 18 lentelės), panaudojimo atvejis PA-18.
    \item \textbf{Aprašymas:} Fono procesas turi automatiškai sugeneruoti ir išsiųsti el. pašto pranešimą aktyvią rezervaciją turinčiam dalyviui apie artėjantį užsiėmimą.
    \item \textbf{Patikrinamas rezultatas:} Nustatytu laiku (pvz., likus 2 val. iki pradžios) dalyvis gauna automatinį el. laišką su užsiėmimo informacija.
\end{itemize}

\subsection{P4. Asmeninių programų valdymas}

\vspace{0.3cm}
\noindent\textbf{FR-36: Dalyvio įvertinimo įrašymas}
\begin{itemize}
    \item \textbf{Užsakovo poreikis ir funkcijos:} F18 (atitinka funkcijos kodą P4.2 iš projekto plano 18 lentelės), panaudojimo atvejis PA-13.
    \item \textbf{Aprašymas:} Treneris turi turėti galimybę įvesti ir išsaugoti dalyvio fizinio įvertinimo ir būklės duomenis.
    \item \textbf{Patikrinamas rezultatas:} Įvertinimo duomenys susiejami su dalyvio kortele ir yra pasiekiami treneriui sudarant programą.
\end{itemize}

\vspace{0.3cm}
\noindent\textbf{FR-37: Asmeninės programos sudarymas ir rekomendacijos}
\begin{itemize}
    \item \textbf{Užsakovo poreikis ir funkcijos:} F18 (atitinka funkcijos kodus P4.1, P4.3 iš projekto plano 18 lentelės), panaudojimo atvejis PA-13.
    \item \textbf{Aprašymas:} Treneris turi galimybę pagal įvertinimą sudaryti dalyviui asmeninę programą bei įrašyti pasiruošimo rekomendacijas.
    \item \textbf{Patikrinamas rezultatas:} Asmeninė programa išsaugoma ir tampa matoma dalyvio paskyroje.
\end{itemize}

\vspace{0.3cm}
\noindent\textbf{FR-38: Namų užduočių priskyrimas}
\begin{itemize}
    \item \textbf{Užsakovo poreikis ir funkcijos:} F18 (atitinka funkcijos kodą P4.6 iš projekto plano 18 lentelės), panaudojimo atvejis PA-14.
    \item \textbf{Aprašymas:} Treneris turi galimybę priskirti dalyviui konkrečias namų užduotis iš pratimų katalogo.
    \item \textbf{Patikrinamas rezultatas:} Priskirta užduotis rodoma dalyvio asmeninių užduočių sąraše.
\end{itemize}

\vspace{0.3cm}
\noindent\textbf{FR-39: Nuotolinis asmeninės programos ir užduočių redagavimas}
\begin{itemize}
    \item \textbf{Užsakovo poreikis ir funkcijos:} F19 (atitinka funkcijos kodus P4.4, P4.7 iš projekto plano 18 lentelės), panaudojimo atvejai PA-13, PA-14.
    \item \textbf{Aprašymas:} Treneris turi turėti galimybę nuotoliniu būdu redaguoti paskirtą asmeninę programą bei pakeisti namų užduotis.
    \item \textbf{Patikrinamas rezultatas:} Redaguoti duomenys išsaugomi, o dalyvis savo paskyroje iškart mato atnaujintą versiją.
\end{itemize}

\vspace{0.3cm}
\noindent\textbf{FR-40: Asmeninės programos archyvavimas ir užduočių šalinimas}
\begin{itemize}
    \item \textbf{Užsakovo poreikis ir funkcijos:} F19 (atitinka funkcijos kodus P4.5, P4.8 iš projekto plano 18 lentelės), panaudojimo atvejai PA-13, PA-14.
    \item \textbf{Aprašymas:} Treneris turi turėti galimybę užbaigti arba suarchyvuoti asmeninę programą bei pašalinti nebematomas namų užduotis.
    \item \textbf{Patikrinamas rezultatas:} Suarchyvuota programa perkeliama į istoriją, o pašalinta užduotis neberodoma aktyvių sąraše.
\end{itemize}

\vspace{0.3cm}
\noindent\textbf{FR-41: Asmeninių programų sąrašo ir detalių peržiūra}
\begin{itemize}
    \item \textbf{Užsakovo poreikis ir funkcijos:} F18, F19 (atitinka funkcijos kodus P4.9, P4.10 iš projekto plano 18 lentelės), panaudojimo atvejai PA-13, PA-14.
    \item \textbf{Aprašymas:} Treneris (visų savo sudarytų) ir dalyvis (savo asmeninės) turi turėti galimybę peržiūrėti asmeninių programų sąrašą bei išsamų jų turinį.
    \item \textbf{Patikrinamas rezultatas:} Sąsajoje atvaizduojami visi programos etapai, pratimai bei pridėtos rekomendacijos.
\end{itemize}

\subsection{P5. Vaizdinio turinio pateikimas}

\vspace{0.3cm}
\noindent\textbf{FR-42: Užsiėmimo etapo vaizdo demonstravimas ekrane}
\begin{itemize}
    \item \textbf{Užsakovo poreikis ir funkcijos:} F9 (atitinka funkcijos kodą P5.1 iš projekto plano 18 lentelės), panaudojimo atvejis PA-15.
    \item \textbf{Aprašymas:} Vietinio grotuvo programa ekrane turi rodyti esamo etapo vaizdo medžiagą, tekstines instrukcijas, atbulinį laiko skaičiavimą bei kito etapo anonsą.
    \item \textbf{Patikrinamas rezultatas:} Salės ekrane sklandžiai atkuriamas vaizdas ir teisingai rodomas laikmatis.
\end{itemize}

\vspace{0.3cm}
\noindent\textbf{FR-43: Autonominis grotuvo veikimas nutrūkus ryšiui}
\begin{itemize}
    \item \textbf{Užsakovo poreikis ir funkcijos:} F8 (atitinka funkcijos kodus P2.14, S2 iš projekto plano), panaudojimo atvejis PA-16.
    \item \textbf{Aprašymas:} Nutrūkus ryšiui tarp grotuvo ir serverio, vietinis grotuvas turi pabaigti pradėtą etapą iš vietinės talpyklos, nepradėti naujo etapo ir informuoti budintį darbuotoją.
    \item \textbf{Patikrinamas rezultatas:} Atjungus tinklą vykdymo metu, grotuvas pabaigia groti esamą įrašą, sustoja ir pateikia pranešimą sąsajoje.
\end{itemize}

\subsection{P6. Garso įranga}

\vspace{0.3cm}
\noindent\textbf{FR-44: Garso instrukcijų ir trenerio balso transliacija}
\begin{itemize}
    \item \textbf{Užsakovo poreikis ir funkcijos:} F10, F11, F16 (atitinka funkcijos kodus P6.1, P6.2 iš projekto plano 18 lentelės), panaudojimo atvejis PA-15.
    \item \textbf{Aprašymas:} Garso subsistema turi perduoti įrašytas garso instrukcijas į ausines/kolonėles bei leisti treneriui transliuoti savo balsą per mikrofoną, automatiškai prislopinant foninį garsą.
    \item \textbf{Patikrinamas rezultatas:} Įrašytos instrukcijos girdimos įrenginiuose, o treneriui prabilus į mikrofoną, muzikos garsas automatiškai nutildomas/sumažinamas.
\end{itemize}

\subsection{P7. Duomenų ir istorijos valdymas}

\vspace{0.3cm}
\noindent\textbf{FR-45: Užsiėmimų eigos istorijos saugojimas}
\begin{itemize}
    \item \textbf{Užsakovo poreikis ir funkcijos:} F13 (atitinka funkcijos kodą P7.1 iš projekto plano 18 lentelės), panaudojimo atvejis PA-17.
    \item \textbf{Aprašymas:} Atsakingas modulis turi automatiškai išsaugoti įvykusių užsiėmimų duomenis, jų faktinę vykdymo trukmę bei susijusius žurnalus.
    \item \textbf{Patikrinamas rezultatas:} Pasibaigus užsiėmimui, jo įrašas atsiranda istorijoje su tiksliais laiko ir etapų parametrais.
\end{itemize}

\vspace{0.3cm}
\noindent\textbf{FR-46: Lankomumo ir užimtumo ataskaitų generavimas}
\begin{itemize}
    \item \textbf{Užsakovo poreikis ir funkcijos:} F13 (atitinka funkcijos kodus P7.2, P7.3 iš projekto plano 18 lentelės), panaudojimo atvejis PA-17.
    \item \textbf{Aprašymas:} Administratorius ir treneris turi turėti galimybę sugeneruoti lankomumo bei erdvių/įrangos užimtumo ataskaitas pasirinktam laikotarpiui (įskaitant virtualaus trenerio užsiėmimus).
    \item \textbf{Patikrinamas rezultatas:} Pasirinkus datų rėžius, sistema sugeneruoja ataskaitą su apskaičiuotais lankomumo skaičiais ir procentiniu užimtumu.
\end{itemize}



\end{document}
